\documentclass[12pt]{article}

\usepackage[T1]{fontenc}
\usepackage[utf8]{inputenc}
\usepackage[margin=1in]{geometry}
\usepackage{setspace}
\usepackage{hyperref}
\usepackage{enumitem}
\usepackage{xcolor}
\usepackage{framed}
\usepackage{graphicx}

\hypersetup{
    colorlinks=true,
    linkcolor=blue,
    urlcolor=blue,
    citecolor=blue
}

\setstretch{1.15}
\setlist[itemize]{topsep=4pt,itemsep=3pt,parsep=2pt}
\setlist[enumerate]{topsep=4pt,itemsep=3pt,parsep=2pt}

\title{Literacy as Power: Bridging Social Justice and Critical Communication\\
\large Communication Access, Neurodivergence, and the Cost of Being Misread in STEM}
\author{Piter Garcia}
\date{Summer 2026}

\begin{document}

\maketitle

\begin{center}
\textit{Draft for EDU498: Exploring Literacy as Power: Bridging Social Justice and Critical Communication}
\end{center}

\section*{Opening: The Bridge}

Literacy is often described as the ability to read and write, but in school it becomes something larger and more dangerous: it becomes the power to be understood. A student must read the syllabus, interpret the tone of an email, decode what counts as an official requirement, understand whether a verbal instruction has the force of a written rule, know how accommodations work, recognize when to contact Disability Services, and decide when a teacher's decision is simply difficult or actually inaccessible. That is not only academic literacy. It is institutional survival literacy.

This journal begins from a documented adult experience in a BIO/STEM course, but the purpose is not to tell a private grievance story. The purpose is to ask what happens when communication itself becomes the barrier. If an adult neurodivergent and chronically ill graduate student, with technical expertise, documentation habits, advisor support, and the ability to advocate, can still be pushed into health strain, relational conflict, and institutional exhaustion, then what happens to children and younger students who cannot defend themselves yet?

That question is the bridge between my experience and social justice. I am not writing only because I was hurt. I am writing because the same classroom patterns that harmed me can quietly teach younger neurodivergent students that they are lazy, difficult, careless, dramatic, or broken, when the real problem is that the communication system was never designed for them to access.

\section*{The Social Justice Issue: Communication Access Injustice in STEM}

The social justice issue I am focusing on is communication access injustice for neurodivergent and chronically ill students in STEM education. By communication access injustice, I mean the harm that happens when a course assumes one default learner: a student who remembers verbal instructions, tolerates ambiguity, processes public correction calmly, writes by hand when expected, manages group conflict while sick, interprets hidden grading rules, and already knows how to activate institutional support.

This issue is not only about accommodations in the narrow legal sense. A student can have accommodations on file and still be excluded if the course communication design makes those accommodations difficult to use. A Flex Plan, a DSO letter, or an advisor conversation cannot fix a classroom culture where important expectations are scattered across verbal comments, emails, submission boxes, peer evaluation routines, and unspoken norms. The problem is not only whether support exists. The problem is whether the student can access the language and process of support in time, without repeatedly exposing private medical information, risking retaliation, or being misread as irresponsible.

This matters historically because schools have long used one version of literacy as the measure of intelligence: fast reading, polished writing, standard language, quiet compliance, memory for verbal instruction, and comfort with institutional procedures. Students who do not fit that profile are often treated as less capable, even when they are thinking deeply. In STEM, this becomes even more intense because the field often treats rigor as rigidity. But rigor without access does not produce excellence. It produces exclusion.

\section*{Research and Documentation: What the Evidence Shows}

For this assignment, I gathered evidence from email records, course transcripts, advisor support records, and EDU498 course materials. The point of the evidence is not to prove that one person had bad intentions. The point is to show how a pattern of communication decisions can create an inaccessible learning environment.

\subsection*{Case 1: A Verbal-Only Rule Became an Access Barrier}

One of the strongest documented examples involves a handwritten peer-evaluation requirement. The email record shows that the instructor stated, after a missing peer-evaluation issue, that ``the instructions given in class were very clear that it needed to be hand-written'' and that a second evaluation would not be accepted if typed.\footnote{Evidence source: Selected Gmail Thread Exports, ``Missing Peer Evaluation / Handwritten Requirement,'' March 19, 2026.}

On the surface, this might sound like a small rule: handwritten versus typed. But for a neurodivergent student, the issue was not handwriting alone. The issue was where the rule lived. If the requirement existed only as a classroom statement and not in the syllabus, written assignment instructions, or transcribed notes, then the student was being held accountable to a literacy event that was not equally accessible.

In a follow-up message to an advisor, I documented why this mattered: because of my ADHD, I transcribed every class session precisely so that verbal-only instructions would not fall through the cracks. I wrote that being penalized based on an undocumented, verbal-only rule was not only academically unfair; it was ``inaccessible by design.''\footnote{Evidence source: Selected Gmail Thread Exports, ``Advisor Follow-Up / Formal Escalation and Ombuds,'' March 27--30, 2026; Megan Lehman Advisor Support Index.}

This case shows literacy as power. The power was not only in the rule. It was in the form of the rule. A written rule can be reread, searched, transcribed, copied into a calendar, discussed with DSO, or clarified with an advisor. A verbal-only rule requires memory, presence, processing speed, and confidence that the student heard it correctly. When a graded requirement depends on that kind of access, the classroom is no longer measuring only responsibility. It is measuring whether the student can survive a communication design built around someone else's brain.

\subsection*{Case 2: Health Disclosure Became Team Responsibility}

Another evidence cluster shows how health disclosure and team responsibility became tangled. In one email, I informed the instructor that I could not attend class because I was experiencing a flare-up of multiple ongoing health conditions and that the medication I needed made effective participation difficult. I also stated that I would catch up with the team and follow up with DSO and advisors if the condition persisted. The instructor replied with well wishes and added, ``make sure your team knows how to proceed.''\footnote{Evidence source: Selected Gmail Thread Exports, ``Health-Related Absence and Team Responsibility,'' March 19, 2026.}

The reply may look kind. That is part of why this type of harm is hard to explain. The problem is not that the sentence was openly cruel. The problem is that, in practice, the communication shifted responsibility back onto a sick student at the exact moment when the student's body and medication were already limiting participation. In group-based STEM work, this matters because a student's absence is not just personal. It becomes relational. If the team already has unclear norms, vague expectations, or dependency on one technically skilled member, then ``make sure your team knows how to proceed'' can intensify the burden.

This is where communication access becomes a social justice issue. A chronically ill student should not have to repeatedly prove illness, manage peer emotions, protect team workflow, interpret accommodation boundaries, and maintain credibility all at once. That is not equal access. It is survival under pressure.

\subsection*{Case 3: Accommodations Became a Limit Instead of a Design Floor}

A later health-related absence shows another layer of the problem. After I reported another health episode consistent with DSO accommodations, the instructor replied by asking whether I had gone over the missed-class limit from the Flex Plan and stated, ``I think it said two classes...correct?''\footnote{Evidence source: Selected Gmail Thread Exports, ``Later Health Absence and Flex Plan Limit,'' April 23, 2026.}

This evidence is important because it shows the difference between accommodations as procedural boundaries and accommodations as access design. A Flex Plan may have rules. But if the communication around the plan feels like a countdown instead of support, then the student experiences the accommodation as another system to fear. For students with chronic illness, the problem is not only missing class. The problem is living with a body that cannot always be scheduled, then having to translate that body into institutional language repeatedly.

In the final paper, I do not need to claim that the instructor intended harm. The documented issue is enough: the communication placed a medically vulnerable student in a position where access had to be continually interpreted, defended, and negotiated.

\subsection*{Case 4: Ambiguous Grading Language Became Policy Literacy}

The email record also shows grading-language confusion around weekly reports. The instructor explained that students only needed to submit ten reports across the semester, but later clarified that he would grade only the first ten submitted.\footnote{Evidence source: Selected Gmail Thread Exports, ``Week 7 Submission / First-Ten Policy,'' February 26--March 1, 2026.} I explained that in prior courses I had understood ``10 out of 12'' structures as meaning a set number of required submissions, often with lower scores dropped or later feedback still available. Later, a related Week 4 report thread records me explaining that the syllabus said ``10 out of the 12 available weeks'' and did not explicitly say ``the first 10 out of 12 reports.''\footnote{Evidence source: Selected Gmail Thread Exports, ``Week 4 Report as Formative Feedback Request,'' April 28--May 5, 2026.}

This case shows that syllabus literacy is not simple. A phrase that looks clear to an instructor may not be clear to students who come from different academic systems, different course cultures, or different access needs. ``10 out of 12'' is not just a number. It carries assumptions about grading, feedback, deadlines, formative learning, and whether later work can replace earlier work.

For neurodivergent students, this matters because ambiguity consumes executive function. The student must not only do the work; they must decode the system behind the work. If the system is not written in plain, stable language, then students who rely on explicit structure are forced to spend their energy guessing what the course means instead of learning the course content.

\subsection*{Case 5: Team Communication Tools Were Not Enough Without Norms}

The evidence pack also documents a search for Slack records and related transcripts. Although a verified Slack export archive was not located, BIOL550 transcripts were copied because they discuss Slack, team communication, professor-related issues, accommodations, peer evaluation, and neurodivergence.\footnote{Evidence source: Slack Export and Transcript Search Log, EDU498 evidence pack.} The search log notes that Slack can be an access tool when it creates durable written records, task memory, and asynchronous participation, but it can also intensify harm when messages are misread, indirect, or used without shared norms.\footnote{Evidence source: Slack Export and Transcript Search Log, assignment use section.}

This is a crucial point for critical communication. A tool does not become accessible simply because it is digital. Slack, Google Docs, AI summaries, transcripts, code files, and shared folders can support neurodivergent access, but only if the class has explicit norms for what each tool means. Which tool controls deadlines? Which tool holds final decisions? Where should students ask questions? How should conflict be handled? What counts as a task assignment? What counts as feedback?

Without those norms, digital tools can become another site of misreading. The student who documents heavily can be treated as overwhelming. The student who asks for clarity can be treated as difficult. The student who uses AI or transcripts to manage memory can be treated as suspicious or excessive. The tool is not the solution unless the communication culture is also redesigned.

\section*{Connecting the Evidence to EDU498 Readings}

Module 4 gives me the language to understand why this is a literacy issue and not only a course conflict. Muhammad's five pursuits frame literacy as identity, skills, intellect, criticality, and joy. In this case, each pursuit was involved. Identity was involved because neurodivergence and chronic illness shaped how I communicated, documented, and asked for clarity. Skills were involved because the course required scientific communication, code, reports, data files, and team coordination. Intellect was involved because my technical ability was needed, but my ways of organizing and communicating were not always recognized as legitimate. Criticality was involved because I had to ask how power operated through syllabus language, email tone, grading rules, and institutional pathways. Joy was involved because learning became fear and defense instead of curiosity and growth.

Moje, Luke, Davies, and Street help explain how literacy tasks position students. A student can be positioned as capable or difficult depending on how the classroom reads their communication. Skerrett's idea of students ``hatching'' in a class also matters here. A good classroom gives students conditions to emerge into stronger versions of themselves. An inaccessible classroom can do the opposite: it can make a student shrink, mask, over-document, lose trust, and become more isolated.

Moje's disciplinary literacy work also matters because STEM is not language-free. Bioinformatics requires reading code, reading data, reading instructions, reading scientific conventions, reading team expectations, and reading institutional signals. If a teacher treats scientific literacy as separate from communication access, students can be excluded even when they understand the science.

Baker-Bell's work on language power also helps me see that no classroom language is neutral. In my case, the issue was not Black Language specifically, but the principle still matters: the language practices that schools treat as normal are tied to power. When a course privileges one communication style, one documentation style, one mode of evaluation, or one way of sounding professional, it can make other literacies appear inferior.

Through my Puzzle Plan, this case connects to access, identity, healing, agency, criticality, joy, and transformation. Through EQUITAS, it asks one central question: does the communication design give students real access, dignity, and multiple ways to demonstrate knowledge, or does it quietly sort students by how well they survive ambiguity?

\section*{Connect Issue to the Present: Universal Design for Learning as a Bridge}

For the present-day initiative, I connect this issue to Universal Design for Learning (UDL), especially the CAST UDL Guidelines 3.0. CAST describes the UDL Guidelines as a tool educators can use across disciplines to make meaningful, challenging learning opportunities accessible to all learners. The guidelines foreground multiple means of engagement, representation, and action/expression, with attention to learner agency, welcoming identities, clarifying language and symbols, using multiple media for communication, organizing information and resources, and challenging exclusionary practices.\footnote{CAST. (2024). \textit{CAST Universal Design for Learning Guidelines version 3.0}. \url{https://udlguidelines.cast.org/}.}

UDL is the right bridge because the evidence in this case is not asking for lower expectations. It is asking for better design. A student should not have to choose between rigor and access. A rigorous STEM course can still provide clear written requirements, plain-language summaries, multiple ways to demonstrate thinking, accessible feedback channels, and private accommodation communication.

\begin{framed}
\noindent \textbf{Screenshot placeholder for final submission:} Insert screenshot of the CAST UDL Guidelines 3.0 page showing the categories for Engagement, Representation, and Action \& Expression. In the caption, explain that UDL is used here as the present-day initiative connecting communication access to inclusive course design.
\end{framed}

The target audience for UDL includes educators, curriculum designers, administrators, researchers, parents, and learning designers. For my journal, the most important audience is STEM faculty and K--12 teachers who may not see communication as part of access. UDL helps make that visible. It shows that access is not only about individual accommodation letters. Access is built into goals, methods, materials, assessments, language, tools, feedback, and opportunities for expression.

However, UDL itself can still face communication barriers. If UDL is presented only in professional jargon, busy teachers may treat it as another abstract framework. If it is presented only as a disability framework, teachers may miss that it benefits all students. If it is not tied to concrete classroom practices, it may not change what happens when a student misses a verbal instruction, needs a written requirement, asks for a different expression mode, or is publicly questioned about a strategy connected to disability.

That is why this journal uses UDL but also pushes beyond it through EQUITAS. UDL gives the design foundation. EQUITAS asks the justice question: who is still being misread, and what must change so access is not dependent on a student's ability to fight?

\section*{Recommendations: Communication as Access Design}

Based on the evidence, I propose the following strategies for STEM instructors, teacher educators, and institutions.

\subsection*{1. Put all graded requirements in one stable written location.}

If a requirement affects a grade, it should not live only in a verbal announcement. It should appear in the syllabus, assignment page, rubric, or a clearly labeled update log. Students should not have to prove that they did not hear something.

\subsection*{2. Use plain-language assignment summaries.}

Each assignment should include a short plain-language section: what to submit, where to submit it, how it will be graded, what format is required, what is optional, what counts as late, and where to ask questions.

\subsection*{3. Label the controlling communication channel.}

Courses often use email, myCourses, Slack, Google Docs, verbal announcements, and assignment boxes at the same time. Instructors should state clearly which channel controls deadlines, requirements, feedback, and policy updates.

\subsection*{4. Treat accommodations as the floor, not the ceiling.}

A DSO letter should not be the only access structure. Instructors should design communication that reduces the need for repeated disclosure. Clear instructions, flexible participation pathways, private correction, and multimodal options help everyone.

\subsection*{5. Provide multimodal ways to demonstrate thinking.}

If the goal is peer feedback, students may be able to demonstrate that through handwriting, typed notes, audio, annotated documents, structured forms, or digital comments. If the goal is scientific understanding, students may show it through code, diagrams, reports, oral explanation, visualizations, or collaborative documentation. The mode should serve the learning goal, not become an unnecessary gate.

\subsection*{6. Keep disability-related clarification private and respectful.}

Public questioning can trigger shame, confusion, or forced disclosure. When a student's access strategy appears unusual, the instructor should clarify privately and adult-to-adult.

\subsection*{7. Build team communication norms before conflict happens.}

Group projects should include explicit norms for Slack, shared documents, task trackers, meetings, deadlines, disagreement, and feedback. Neurodivergent access requires more than a tool. It requires agreed-upon meanings for how the tool will be used.

\subsection*{8. Train STEM faculty in accessible rigor.}

Faculty should be trained to understand that rigor is not the same as ambiguity, speed, handwriting, memory, or emotional toughness. Rigor means meaningful challenge with accessible pathways to participation.

\subsection*{9. Create an EQUITAS communication checklist.}

My next step is to turn this experience into an EQUITAS-style communication checklist for teachers. It would ask:

\begin{itemize}
    \item Is every graded expectation written, findable, and stable?
    \item Are verbal updates repeated in writing?
    \item Are students told which channel controls policy and deadlines?
    \item Are accessibility concerns handled privately?
    \item Does the assignment allow multiple ways to demonstrate knowledge when the mode is not the learning goal?
    \item Are peer-team norms explicit before conflict begins?
    \item Does the communication design reduce disclosure burden?
    \item Would this design work for a child who cannot advocate like an adult?
\end{itemize}

This is where EQUITAS becomes more than a reflection framework. It becomes an equivalence framework. Equivalence does not mean every student receives the exact same communication. It means every student receives communication that gives them an equivalent chance to understand, participate, and show what they know.

\section*{Critical Reflection: From Pain to Protection}

What happened in this BIO/STEM course hurt because it did not only affect a grade. It affected my health, my relationships, my trust, and my sense of whether my ways of surviving as a neurodivergent and chronically ill student would ever be read as legitimate. Documentation became survival. Transcripts became memory. Emails became evidence. Advising became translation. Literacy became self-defense.

That is why this assignment matters to me. I do not want to use literacy only to explain what happened after the harm. I want to use literacy to prevent harm before it happens to someone else.

The hardest part is knowing that I had resources many students do not have. I could write carefully. I could document. I could email advisors. I could identify patterns. I could connect what happened to access, institutional literacy, and critical communication. I could survive the experience long enough to turn it into evidence. Many children cannot. Many young students do not have the language to say, ``This is inaccessible by design.'' They may only know that they are overwhelmed, ashamed, angry, sick, or afraid. They may be punished before anyone teaches them how to name the system.

That is why the bridge matters. The goal is not to prove that one professor is the whole problem. The goal is to show that when teachers communicate through hidden rules, single-mode expectations, public pressure, and rigid assumptions, they decide which students are readable as capable. For neurodivergent and chronically ill students, that decision can become physical, emotional, academic, and relational harm.

My Puzzle Plan is about building systems where students are not forced to break themselves to be seen. EQUITAS asks whether learning environments provide real access, dignity, and multiple ways to demonstrate knowledge. This journal adds one more piece: communication itself must be designed as access. If the language of a course cannot be accessed, then the learning cannot be accessed either.

Literacy can exclude, but literacy can also protect. It can document harm. It can translate pain into evidence. It can turn evidence into design. It can turn design into justice. And if we do that work well, maybe the next student, especially the next child, will not have to become their own lawyer, doctor, translator, and advocate just to be treated as a learner.

\section*{Working References and Evidence Sources}

\begin{itemize}
    \item CAST. (2024). \textit{CAST Universal Design for Learning Guidelines version 3.0}. \url{https://udlguidelines.cast.org/}
    \item Garcia, P. (2026). \textit{BIO Professor Communication Evidence Map}. EDU498 evidence file.
    \item Garcia, P. (2026). \textit{Neurodivergence Injustice Evidence Pack}. EDU498 evidence folder.
    \item Garcia, P. (2026). \textit{Selected Gmail Thread Exports}. EDU498 evidence file.
    \item Garcia, P. (2026). \textit{Megan Lehman Advisor Support Index}. EDU498 evidence file.
    \item Garcia, P. (2026). \textit{Slack Export and Transcript Search Log}. EDU498 evidence file.
    \item Garcia, P. (2026). \textit{Module 4 Comprehensive Themed Overview}. EDU498 course synthesis.
\end{itemize}

\end{document}